\documentclass[11pt,a4paper]{article}
\usepackage[utf8]{inputenc}
\usepackage[vietnamese]{babel}
\usepackage{amsmath,amssymb}
\usepackage[version=4]{mhchem}
\usepackage{geometry}
\geometry{top=1.5cm,bottom=1.2cm,left=1.4cm,right=1.4cm,headheight=28pt,headsep=12pt,footskip=18pt}
\usepackage{graphicx}
\usepackage{xcolor}
\usepackage{tabularx}
\usepackage{array}
\usepackage{tikz}
\usepackage{fancyhdr}
\usepackage{enumitem}
\usepackage{multido}
\usepackage{tcolorbox}
\usepackage{colortbl}

\definecolor{hfRed}{HTML}{D32F2F}
\definecolor{hfGreen}{HTML}{8BC390}
\definecolor{mainGreen}{HTML}{2E7D32}
\definecolor{solBlue}{HTML}{1565C0}
\definecolor{tableGreen}{HTML}{E8F5E9}

\usepackage{ex_test}

\pagestyle{fancy}
\fancyhf{}
\renewcommand{\headrulewidth}{0.75pt}
\renewcommand{\headrule}{\hbox to\headwidth{\color{hfGreen}\leaders\hrule height \headrulewidth\hfill}}
\renewcommand{\footrulewidth}{0.75pt}
\renewcommand{\footrule}{\hbox to\headwidth{\color{hfGreen}\leaders\hrule height \footrulewidth\hfill}}

\fancyhead[L]{\small\itshape\color{hfRed}Tài liệu lưu hành nội bộ}
\fancyhead[R]{\small\color{hfRed}Thầy \textbf{TRẦN VĂN THẠNH} - 0777.470.803}

\fancyfoot[L]{\footnotesize\itshape\color{hfRed}CS1: 6/15 Nguyễn Hoàng, P. Kim Long, TP Huế \quad|\quad CS2: 24 Đặng Thái Thân, TP Huế}
\fancyfoot[R]{\small\textbf{\color{hfRed}Trang \thepage}}

\setlength{\parindent}{0pt}
\setlength{\parskip}{2pt}

\begin{document}

\begin{minipage}[t]{0.48\textwidth}
\textbf{TRƯỜNG ĐẠI HỌC KHOA HỌC}\\
\textbf{TRƯỜNG THPT CHUYÊN KHOA HỌC HUẾ}\\
\textit{ĐỀ CHÍNH THỨC}\\
\textit{(Đề thi có 03 trang)}
\end{minipage}
\hfill
\begin{minipage}[t]{0.48\textwidth}
\raggedleft
\textbf{\color{hfRed}ĐỀ KIỂM TRA GIỮA HỌC KỲ I - NĂM HỌC 2026 - 2027}\\
\textbf{Môn: Hóa học -- Lớp 11 (CHUNG)}\\
\textit{Thời gian làm bài: 45 phút (không kể thời gian giao đề)}\\
\textbf{Mã đề: 135}
\end{minipage}

\vspace{4pt}
\noindent
\textit{Họ và tên học sinh:} \dotfill \textit{Lớp:} \dotfill \textit{Số báo danh:} \dotfill

\vspace{3pt}
\noindent
\textit{\footnotesize Cho nguyên tử khối: $\ce{H}=1; \ce{Li}=7; \ce{Be}=9; \ce{N}=14; \ce{O}=16; \ce{Na}=23; \ce{Mg}=24; \ce{Al}=27; \ce{P}=31; \ce{S}=32; \ce{Cl}=35,5; \ce{K}=39; \ce{Ca}=40; \ce{Ba}=137$.}

\vspace{3pt}
\hrule color{hfGreen} height 0.5pt

\vspace{4pt}
\noindent
{\color{mainGreen}\textbf{\Large A. TRẮC NGHIỆM (7 điểm)}}

\vspace{3pt}
\noindent
{\color{mainGreen}\textbf{\large PHẦN I. (3 điểm) Câu trắc nghiệm nhiều phương án lựa chọn.}}\\
\textit{\small Học sinh trả lời từ câu 1 đến câu 12. Mỗi câu hỏi học sinh chỉ chọn một phương án.}
\vspace{2pt}

\begin{ex}
Có thể nhận biết muối ammonium bằng cách cho muối tác dụng với dung dịch kiềm thấy thoát ra một chất khí làm xanh quỳ tím ẩm. Chất khí đó là
\choice
{$\ce{NO2}$}
{$\ce{NO}$}
{\True $\ce{NH3}$}
{$\ce{H2}$}
\end{ex}
\begin{ex}
Trong các dung dịch có cùng nồng độ $0,10\text{ mol/L}$ dưới đây, dung dịch nào dẫn điện kém nhất?
\choice
{\True $\ce{HF}$}
{$\ce{HI}$}
{$\ce{HCl}$}
{$\ce{HBr}$}
\end{ex}
\begin{ex}
Acid trong dạ dày giúp kích hoạt enzyme pepsin, hỗ trợ tiêu hóa protein và bảo vệ cơ thể khỏi vi khuẩn có hại. Nếu trong dạ dày có $\text{pH} = 2$ thì nồng độ $\ce{H+}$ bằng bao nhiêu?
\choice
{$10\text{ M}$}
{$10^2\text{ M}$}
{\True $10^{-2}\text{ M}$}
{$2\text{ M}$}
\end{ex}
\begin{ex}
Nhận xét nào dưới đây về nitrogen \textbf{không đúng}?
\choice
{\True Khí nitrogen duy trì sự cháy}
{$\ce{N2}$ phản ứng với $\ce{O2}$ ở nhiệt độ rất cao}
{Khí nitrogen rất ít tan trong nước}
{$\ce{N2}$ vừa có tính oxi hóa vừa có tính khử}
\end{ex}
\begin{ex}
Phản ứng thuận nghịch là phản ứng
\choice
{có phương trình hoá học được biểu diễn bằng mũi tên một chiều}
{chỉ xảy ra theo một chiều nhất định}
{\True trong cùng điều kiện, phản ứng xảy ra theo hai chiều trái ngược nhau}
{xảy ra giữa hai chất khí}
\end{ex}
\begin{ex}
Các yếu tố ảnh hưởng đến cân bằng hoá học là :
\choice
{nồng độ, nhiệt độ và chất xúc tác}
{nồng độ, áp suất và diện tích bề mặt}
{\True nồng độ, nhiệt độ và áp suất}
{áp suất, nhiệt độ và chất xúc tác}
\end{ex}
\begin{ex}
Hành động nào của con người \textbf{không phải} nguyên nhân chính gây nên hiện tượng phú dưỡng?
\choice
{Sử dụng dư thừa lượng phân bón chứa ion $\ce{PO4^3-}$}
{\True Quá trình đánh bắt cá bằng chích điện và thuốc nổ của con người}
{Sử dụng dư thừa lượng phân bón chứa ion $\ce{NO2^-}, \ce{NO3^-}$}
{Nước thải chăn nuôi, nước thải công nghiệp chứa nhiều các nguyên tố dinh dưỡng như N, P}
\end{ex}
\begin{ex}
Trong phản ứng sau đây, những chất, ion nào đóng vai trò là acid theo thuyết Bronsted-Lowry?
$$\ce{H2S(aq) + H2O <=> HS^-(aq) + H3O^+(aq)}$$
\choice
{$\ce{H2S}$ và $\ce{HS^-}$}
{$\ce{H2O}$ và $\ce{H3O^+}$}
{$\ce{H2S}$ và $\ce{H2O}$}
{\True $\ce{H2S}$ và $\ce{H3O^+}$}
\end{ex}
\begin{ex}
Nitrogen tương đối trơ về mặt hoá học ở nhiệt độ thường là do
\choice
{\True phân tử $\ce{N2}$ có liên kết ba với năng lượng liên kết lớn}
{nitrogen có độ âm điện lớn}
{phân tử $\ce{N2}$ có liên kết cộng hoá trị không phân cực}
{phân tử $\ce{N2}$ có liên kết ion}
\end{ex}
\begin{ex}
Khi tăng áp suất, những cân bằng hóa học nào dưới đây chuyển dịch theo chiều thuận?\\
(1) $\ce{H2(g) + Cl2(g) <=> 2HCl(g)}$\\
(2) $\ce{NH4Cl(s) <=> NH3(g) + HCl(g)}$\\
(3) $\ce{2H2O2(aq) <=> 2H2O(l) + O2(g)}$\\
(4) $\ce{3O2(g) <=> 2O3(g)}$
\choice
{(1) và (2)}
{(1) và (4)}
{(2) và (3)}
{\True Chỉ có (4)}
\end{ex}
\begin{ex}
Dãy chất nào sau đây đều là chất điện li yếu?
\choice
{$\ce{H2S}, \ce{H2SO3}, \ce{H2SO4}$}
{$\ce{H2CO3}, \ce{H3PO4}, \ce{CH3COOH}, \ce{Ba(OH)2}$}
{\True $\ce{H2S}, \ce{CH3COOH}, \ce{HClO}$}
{$\ce{H2CO3}, \ce{H2SO3}, \ce{HClO}, \ce{Al2(SO4)3}$}
\end{ex}
\begin{ex}
Dung dịch nào sau đây có $\text{pH} > 7$?
\choice
{\True $\ce{NaOH}$}
{$\ce{H2SO4}$}
{$\ce{NaCl}$}
{$\ce{HNO3}$}
\end{ex}
\vspace{4pt}
\noindent
{\color{mainGreen}\textbf{\Large PHẦN II (2 điểm). Câu trắc nghiệm đúng sai.}}\\
\textit{\small Học sinh trả lời từ câu 1 đến câu 2. Trong mỗi ý a), b), c), d) ở mỗi câu, học sinh chọn đúng hoặc sai.}
\vspace{2pt}

\begin{ex}
Trong khí thải của quy trình sản xuất thuốc trừ sâu, phân bón hóa học có lẫn khí $\ce{NH3}$. Khí này rất độc đối với sức khoẻ của con người và gây ô nhiễm môi trường. Con người hít phải khí này với lượng lớn sẽ gây ngộ độc: ho, đau ngực (nặng: đau thắt ngực, khó thở, thở nhanh, thở khò khè; chảy nước mắt và bỏng mắt, mù mắt, đau họng nặng, đau miệng; mạch nhanh, yếu, sốc; lẫn lộn, đi lại khó khăn, chóng mặt, thiếu sự phối hợp, bồn chồn, ngẩn ngơ).
\choiceTF
{\True Để xử lí $\ce{NH3}$ lẫn trong khí thải, người ta có thể dẫn khí thải qua một bể lọc chứa dung dịch $\ce{HCl}$}
{Tất cả các muối của $\ce{NH3}$ đều bền với nhiệt}
{$\ce{NH3}$ có tính base mạnh và tính khử mạnh}
{\True Đưa hai đũa thuỷ tinh có nhúng $\ce{NH3}$ đặc và $\ce{HCl}$ đặc lại gần thấy khói màu trắng xuất hiện}
\end{ex}
\begin{ex}
Trong bình kín có hệ cân bằng hóa học sau:
$$\ce{CO2(g) + H2(g) <=> CO(g) + H2O(g)};\quad \Delta_{\mathrm{r}}H^{\circ}_{298} > 0$$
Cho các phát biểu về sự tác động đến hệ cân bằng dưới đây, hãy xét tính đúng sai của các phát biểu đó.
\choiceTF
{\True Tăng nhiệt độ của hệ thì cân bằng chuyển dịch theo chiều thuận}
{Thêm chất xúc tác thì cân bằng chuyển dịch theo chiều thuận}
{Giảm áp suất chung của hệ thì cân bằng chuyển dịch theo chiều nghịch}
{Nếu nồng độ ban đầu $\ce{CO2}$ và $\ce{H2}$ đều là $0,03\text{ M}$; tại cân bằng $[\ce{CO}] = 0,018\text{ M}$ thì giá trị $K_c$ là $1,5$}
\end{ex}
\vspace{4pt}
\noindent
{\color{mainGreen}\textbf{\Large PHẦN III (2 điểm). Câu trắc nghiệm trả lời ngắn.}}\\
\textit{\small Học sinh trả lời từ câu 1 đến câu 4.}
\vspace{2pt}

\begin{ex}
Trong công nghiệp ammonia được tổng hợp từ nitrogen và hydrogen. Cho $14,874\text{ lít }\ce{N2}$ (đkc) tác dụng với lượng dư khí $\ce{H2}$. Biết hiệu suất của phản ứng là $30\%$, khối lượng $\ce{NH3}$ tạo thành là bao nhiêu gam? \textit{(Kết quả làm tròn đến hàng phần trăm)}

\noindent\textit{Trả lời:} \dotfill
\end{ex}
\begin{ex}
Cho các chất dưới đây: $\ce{HClO4}, \ce{HClO}, \ce{HF}, \ce{HNO3}, \ce{H2S}, \ce{H2SO3}, \ce{NaOH}, \ce{NaCl}, \ce{CuSO4}, \ce{CH3COOH}$. Số chất thuộc loại chất điện li mạnh là bao nhiêu?

\noindent\textit{Trả lời:} \dotfill
\end{ex}
\begin{ex}
Thực hiện thí nghiệm chuẩn độ dung dịch $\ce{HCl}$ chưa biết nồng độ bằng dung dịch chuẩn $\ce{NaOH}\ 0,5\text{M}$ và thuốc thử phenolphthalein. Cho $10,00\text{ mL}$ dung dịch $\ce{HCl}$ vào bình tam giác cùng với 1 -- 2 giọt thuốc thử. Thể tích dung dịch $\ce{NaOH}$ ở các lần chuẩn độ được ghi lại như sau:
\begin{center}
\renewcommand{\arraystretch}{1.2}
\begin{tabular}{|c|c|c|}
\hline
\rowcolor{tableGreen}
\textbf{Lần 1} & \textbf{Lần 2} & \textbf{Lần 3} \\
\hline
$15,1\text{ mL}$ & $14,7\text{ mL}$ & $15,2\text{ mL}$ \\
\hline
\end{tabular}
\end{center}
Dung dịch $\ce{HCl}$ có nồng độ bằng bao nhiêu mol/L? \textit{(Kết quả làm tròn đến hàng phần trăm)}

\noindent\textit{Trả lời:} \dotfill
\end{ex}
\begin{ex}
Cho phản ứng: $\ce{CO(g) + H2O(g) <=> CO2(g) + H2(g)}\quad \Delta_{\mathrm{r}}H^{\circ}_{298} = -41\text{ kJ}$.\\
Các phát biểu sau đây, có bao nhiêu phát biểu \textbf{đúng} về phản ứng trên?\\
(1) Phản ứng theo chiều thuận là phản ứng thu nhiệt.\\
(2) Tăng nhiệt độ cân bằng chuyển dịch theo chiều nghịch.\\
(3) Thêm khí $\ce{H2}$ vào hệ cân bằng chuyển dịch theo chiều nghịch.\\
(4) Tăng áp suất chung bằng cách nén cho thể tích của hệ giảm xuống cân bằng chuyển dịch theo chiều thuận.

\noindent\textit{Trả lời:} \dotfill
\end{ex}
\vspace{6pt}
\noindent
{\color{mainGreen}\textbf{\Large B. TỰ LUẬN (3 điểm)}}
\vspace{2pt}

\noindent
\textbf{Câu 1. (1,0 điểm)} Viết phương trình điện li của các chất sau trong nước: $\ce{KOH}$, $\ce{Al2(SO4)3}$, $\ce{CH3COOH}$, $\ce{(NH4)3PO4}$.

\vspace{1pt}
\noindent
\multido{\n=1+1}{6}{\dotfill\\[4pt]}

\vspace{4pt}
\noindent
\textbf{Câu 2. (1,0 điểm)} Trong công nghiệp, người ta sản xuất nitric acid ($\ce{HNO3}$) từ ammonia theo sơ đồ chuyển hoá sau:
$$\ce{NH3 ->[+O2, t^\circ, xt] NO ->[+O2] NO2 ->[+O2 + H2O] HNO3}$$
a) Viết các phương trình hoá học xảy ra.\\
b) Để điều chế $200\,000\text{ tấn}$ nitric acid có nồng độ $60\%$ cần dùng bao nhiêu tấn ammonia? Biết rằng hiệu suất của quá trình sản xuất nitric acid theo sơ đồ trên là $96,2\%$.

\vspace{1pt}
\noindent
\multido{\n=1+1}{8}{\dotfill\\[4pt]}

\vspace{4pt}
\noindent
\textbf{Câu 3. (1,0 điểm)} Trộn $80\text{ mL}$ dung dịch $\ce{HCl}\ 0,5\text{M}$ với $120\text{ mL}$ dung dịch $\ce{NaOH}\ 0,6\text{M}$. Tính pH của dung dịch thu được. \textit{(Kết quả làm tròn đến hàng phần mười)}.

\vspace{1pt}
\noindent
\multido{\n=1+1}{6}{\dotfill\\[4pt]}

\vspace{6pt}
\begin{center}
\textbf{\large ---------- HẾT ----------}\\
\textit{\small - Học sinh không được sử dụng tài liệu; Giám thị không giải thích gì thêm. -}
\end{center}

\end{document}
